% TOP_PROBLEM: {"id": 2644, "title": "Kourovka Notebook Problem 21.135", "category_id": 17, "category": {"id": 17, "name": "group_theory", "display_name": "Group Theory", "description": "Problems about groups, group actions, representations, and related algebraic structures.", "slug": "group-theory", "order_index": 17, "created_at": "2026-07-31T15:26:25.671Z"}, "status": "open", "problem_number": "KOU-21.135", "set_id": 10}
% FIRST_POSTED: 2026-10-01
% ENTRY_KIND: statement_only
% UnsolvedMath ID 2644; source catalogue code KOU-21.135
% !TeX program = lualatex
% Definition and sources only; this file is not an LLM solution attempt.
\documentclass[11pt,a4paper]{article}
\usepackage[margin=25mm]{geometry}
\usepackage{fontspec}
\setmainfont{lmroman10-regular.otf}[BoldFont=lmroman10-bold.otf,
 ItalicFont=lmroman10-italic.otf,BoldItalicFont=lmroman10-bolditalic.otf]
\usepackage{amsmath,amssymb,mathtools}
\usepackage{unicode-math}
\setmathfont{latinmodern-math.otf}
\AtBeginDocument{\let\setminus\smallsetminus}
\usepackage{xurl,hyperref}
\hypersetup{colorlinks=true,urlcolor=blue}
\setcounter{secnumdepth}{0}
\setlength{\parindent}{0pt}
\setlength{\parskip}{5pt}
\setlength{\emergencystretch}{3em}
\begin{document}
\section*{Kourovka Notebook Problem 21.135}\label{2644}
{\small UnsolvedMath ID 2644 · KOU-21.135 · Group Theory · Kourovka Notebook - New Problems, Issue 21}
\subsection{Definitions and mathematical statement}
Let $G$ be a finite group. A complex representation is a homomorphism $\rho:G\to\operatorname{GL}(V)$ with $V$ a finite-dimensional complex vector space; it is irreducible if $V\ne0$ has no proper nonzero invariant subspace. Its character is $\chi(g)=\operatorname{tr}(\rho(g))$, and its degree is $\chi(1)=\dim V$.

Form the multiset
\[
 D(G)=\{\!\{\chi(1):\chi\in\operatorname{Irr}(G)\}\!\},
\]
where each irreducible character is included once. Thus if several inequivalent irreducible representations have the same dimension, that dimension occurs repeatedly. Equivalently, record for each integer $d\ge1$ the number of irreducible representations of dimension $d$.

For a finite group $X$, its solvable radical $R(X)$ is its largest solvable normal subgroup. Solvability means that the derived series $X^{(0)}=X$, $X^{(i+1)}=[X^{(i)},X^{(i)}]$ eventually reaches $1$.

Suppose $G,H$ are finite groups with $D(G)=D(H)$ and $R(G)=1$. Must $R(H)=1$? This is precisely equality of the source's degree data $\chi_1(G)$ including all multiplicities. The familiar identity $|G|=\sum_{\chi\in\operatorname{Irr}(G)}\chi(1)^2$ shows that the hypothesis also fixes the group order. It does not supply character values on individual conjugacy classes. The desired conclusion concerns the solvable radical alone and does not assert that $G$ and $H$ are isomorphic.
\subsection{Short English statement}
Does the complete multiset of irreducible complex character degrees determine whether a finite group has trivial solvable radical?
\subsection{Sources}
\begin{itemize}
\item[S1] ULAM.AI and the UnsolvedMath Contributors, supplied problems.json, record 2644 (KOU-21.135), Kourovka Notebook - New Problems, Issue 21. Catalogue identity and source transcription; CC BY 4.0.
\url{https://huggingface.co/datasets/ulamai/UnsolvedMath}.
\item[S2] The Kourovka Notebook, 21st edition, version 47 (30 September 2026), new problems, Problem 21.135. Expanded formulation of the supplied catalogue statement.
\url{https://arxiv.org/pdf/1401.0300v47}.
\end{itemize}
\end{document}
